\chapter[Considerações Finais]{Considerações Finais}
\label{cap:consideracoes}

\section{Síntese dos resultados e atendimento aos objetivos}

O trabalho tratou da seleção e da distribuição de atrações turísticas em agendas diárias sujeitas a horários, durações e deslocamentos. O desenvolvimento do Interlude reuniu essas decisões em um aplicativo móvel com persistência e edição de viagens, ampliando o fluxo da prova de conceito \textit{web} do TCC1.

O objetivo de analisar a literatura foi desenvolvido na Seção~\ref{cap:teoria}, ao relacionar o planejamento turístico aos problemas de orientação e roteamento e examinar recursos de ferramentas comerciais. Essa base ajudou a delimitar o problema e os requisitos da Seção~\ref{cap:requisitos}. A modelagem apresentada na Seção~\ref{sec:modelagem-matematica} traduziu as escolhas do viajante em restrições e penalidades, explicitando como horários, deslocamentos e prioridades participam da distribuição das visitas.

Os objetivos de projetar a arquitetura e implementar o fluxo móvel foram concretizados na separação entre aplicativo, orquestração, persistência e otimização, descrita na Seção~\ref{cap:solucao}. Essa organização permite preparar dados de provedores e estimativas de modelos de linguagem antes de construir as rotas. No aplicativo, o resultado pode ser salvo e reaberto. Mudanças de datas geram uma nova distribuição, enquanto alterações de dia passam por ajuste local ou por uma proposta de reorganização examinada pelo viajante antes de sua aplicação.

A verificação funcional reuniu oito propriedades em cinco cenários sintéticos e três execuções da integração. Na execução detalhada de Nova Iorque, uma instância de 52 nós distribuiu as doze atrações selecionadas em cinco dias, com 5,032 segundos de busca. Os dias ocuparam entre 2,8 e 8,9 horas, e o deslocamento total foi de 62,4~km. Brasília e Pirenópolis complementaram a observação no ambiente publicado e em aparelho físico. Cada registro contribui com uma perspectiva do funcionamento, conforme os diferentes intervalos de medição apresentados na Seção~\ref{sec:resultados}.

O experimento computacional acrescentou uma comparação entre estratégias sob a mesma entrada. A GLS alcançou objetivo 154.192 nos limites testados entre 2 e 600 segundos, enquanto a descida gulosa obteve 155.336 em aproximadamente 0,11 segundo no registro original. O resultado justifica investigar o compromisso entre tempo e objetivo em outras entradas, mas não uma redução geral do prazo de produção. Os recortes de três e dois dias também mostraram os efeitos das penalidades: pausas automáticas foram omitidas, dias ultrapassaram nove horas e, no período mais curto, quatro atrações ficaram de fora.

A avaliação com seis participantes acrescentou a perspectiva de uso, com 30 casos de teste em seis aparelhos e sete cenários de geração. Os roteiros desses cenários foram resolvidos pelo \textit{solver}, que respeitou os fechamentos semanais e excluiu apenas as últimas prioridades quando a capacidade foi insuficiente. Dos casos de teste, 16 atenderam à expectativa, 9 tiveram ressalvas e 5 não atenderam. Os problemas mais relevantes foram uma reserva de restaurante agendada às 13h56 em vez de 12h30, a ordenação apenas por arraste, que exigiu ajuda no caso CT08, e o dia da chegada com saída da hospedagem antes da chegada. A análise técnica dos cenários mostrou ainda que o piso de visitas não é cobrado de um dia vazio. As respostas coletadas registram a percepção dos participantes e orientam correções, dentro do alcance dessa avaliação.

Essas entregas atendem ao objetivo de desenvolver e verificar funcionalmente a proposta nos cenários examinados. A contribuição está em tornar as escolhas do viajante parte de um processo que obtém informações, organiza as visitas e apresenta uma agenda passível de revisão. A ampliação da avaliação com viajantes permitirá examinar a adequação dessa organização a outras formas de viajar.

\section{Análise crítica e limitações}

Uma das dificuldades da implementação foi transformar informações de fontes distintas em valores que o modelo pudesse utilizar. Quando falta uma rota entre dois lugares, a preparação pode recorrer a uma estimativa geodésica. Quando falta a duração de uma visita, o modelo de linguagem fornece uma estimativa complementar. Essas alternativas permitem continuar o planejamento, mas o erro de um valor pode repercutir nas atividades seguintes. Uma visita que dure mais do que o previsto, por exemplo, pode exigir rever a agenda mesmo que todos os horários tenham sido respeitados no cálculo.

A edição tornou essa dependência entre atividades particularmente visível. Transferir uma atração para outro dia pode afetar o horário de uma visita que o usuário não pretendia alterar. Por isso, o aplicativo distingue um ajuste que preserva o restante da agenda de uma reorganização mais ampla e permite examinar a proposta antes de substituir a viagem salva. A solução trata o conflito como parte do planejamento. Ampliar a observação desse fluxo com outros viajantes ajudará a avaliar se as alternativas e seus efeitos são compreendidos.

A análise dos resultados considera também as condições em que o trabalho foi desenvolvido e verificado. Os pontos abaixo delimitam as conclusões e indicam quais evidências ainda precisam ser produzidas:
\begin{itemize}
  \item \textbf{Experiência dos autores e alcance da avaliação.} Os autores contribuíram com sua experiência de planejamento e com observações obtidas durante os percursos de verificação. Como os autores também desenvolveram o aplicativo, essas observações estão vinculadas à sua familiaridade com o sistema. A avaliação com seis participantes acrescenta dados coletados sobre a execução das tarefas e a percepção de uso. A quantidade de participantes, as diferenças entre os cenários e o emprego de um formulário próprio limitam a generalização dos achados.
  \item \textbf{Cobertura da verificação.} A suíte sintética examina propriedades determinadas, e os testes dirigidos acompanham partes da edição, persistência, calendário e reutilização de dados. As três execuções integradas mostram o funcionamento em condições específicas. Ampliar as entradas e repetir os percursos permitirá investigar comportamentos ainda não cobertos e caracterizar a variação do desempenho.
  \item \textbf{Qualidade da otimização.} A busca trabalha com um prazo e devolve a solução encontrada, sem um certificado de otimalidade. A melhor solução observada no experimento é uma referência de comparação entre execuções, sem garantia de ótimo. Como os pesos foram ajustados durante o desenvolvimento, uma análise de sensibilidade ajudará a compreender seu efeito em outros cenários. O comportamento de exclusão foi verificado em sobrecarga sintética e em um recorte controlado da entrada de Nova Iorque. As execuções integradas de Nova Iorque e Brasília alocaram todas as atrações selecionadas. A comparação permanece restrita a uma entrada e suas derivações, sem amostra de cidades ou variação sistemática dos pesos.
  \item \textbf{Precisão dos dados.} A confiança declarada pelo modelo de linguagem ainda precisa ser confrontada com erros observados, assim como as durações e recomendações de sazonalidade precisam de aferição sistemática. Os deslocamentos também contêm aproximações: a matriz não considera a partida específica de cada visita, e o cálculo geodésico não confirma um caminho praticável nem fornece um limite superior para sua duração. Essas condições devem acompanhar a interpretação dos horários gerados.
  \item \textbf{Dependências operacionais.} Latência e disponibilidade dependem de Places, Routes e Gemini. A reutilização de respostas reduz consultas repetidas, mas pode manter dados desatualizados. O prazo de 90 dias da matriz é uma configuração que ainda precisa ser adequada às permissões de armazenamento do provedor, como discutido na Seção~\ref{sec:custos-geracao}. Os registros de execução ainda não compõem uma série histórica integrada de custo, tempo e características das instâncias.
  \item \textbf{Delimitação do modelo.} O orçamento monetário e a otimização do transporte entre cidades não fazem parte da função objetivo. Os horários logísticos são informados pelo viajante, mas sua tradução para janelas das atrações e das bases precisa ser verificada em conjunto. O registro de Nova Iorque mostrou uma diferença entre a disponibilidade de visitas e o horário permitido de saída da base, discutida na Seção~\ref{sec:execucao-ny}. A reorganização depende de solicitação explícita.
\end{itemize}

\section{Continuidade do estudo}

Esta versão reúne a implementação, a verificação funcional, o experimento computacional e a avaliação realizada com seis participantes. Os resultados orientam melhorias no aplicativo e no modelo de geração de roteiros. As propostas de continuidade são:
\begin{enumerate}
  \item Priorizar as correções da Tabela~\ref{tab:participantes-problemas} que afetam a agenda e a conclusão das tarefas: usar o horário reservado como início da refeição, oferecer alternativa à ordenação por arraste e limitar a saída da hospedagem pelos horários logísticos. Avaliar também se um dia sem visitas deve receber o custo do piso.
  \item Executar cenários repetíveis com diferentes cidades, densidades de atrações e ritmos, separando tempos de preparação, busca e apresentação. Incluir sobrecarga para avaliar os pontos excluídos e a distribuição entre dias.
  \item Ampliar a comparação inicial entre estratégias para outras entradas e instâncias de TOPTW da literatura \cite{vansteenwegen2011survey}, explicitando adaptações necessárias e diferenças entre funções objetivo antes de calcular desvios em relação a soluções conhecidas.
  \item Medir a sensibilidade dos pesos da curadoria e das penalidades, mantendo constantes os demais dados de cada instância, para identificar efeitos sobre deslocamento, exclusões e distribuição temporal.
  \item Aferir durações e recomendações estimadas contra dados de referência e comparar Gemini com um segundo modelo sob as mesmas entradas. Registrar divergências, respostas sem suporte, latência e consumo. A concordância entre modelos será analisada separadamente do acerto, pois não substitui a conferência com fontes externas. O procedimento preparado está em \path{tcc/fontes/avaliacao-llm/}.
  \item Instrumentar o histórico de latência, uso de dados armazenados e custos externos. Avaliar o efeito da atualização da matriz e dos prazos de expiração sobre os roteiros.
\end{enumerate}

O Interlude implementa um fluxo de planejamento no qual prioridades, horários e deslocamentos participam da geração e da revisão de roteiros. Os testes, o experimento computacional e a avaliação com participantes documentam seu comportamento nos cenários examinados. As propostas de continuidade concentram-se na correção dos problemas observados, na aferição dos dados e na ampliação das comparações entre entradas e estratégias de otimização.
